\documentclass[10pt]{article}

\usepackage[letterpaper,margin=0.58in]{geometry}
\usepackage{amsmath,amssymb}
\usepackage{enumitem}
\usepackage{xcolor}
\usepackage{fancyhdr}

\setlength{\parindent}{0pt}
\setlength{\parskip}{0pt}

\pagestyle{fancy}
\fancyhf{}
\renewcommand{\headrulewidth}{0pt}
\cfoot{\thepage}

\definecolor{sectiongray}{gray}{0.935}

\newcommand{\sectionbar}[1]{%
  \vspace{0.55em}%
  \noindent\colorbox{sectiongray}{%
    \parbox{\dimexpr\textwidth-2\fboxsep\relax}{%
      \textbf{\large\strut #1\strut}%
    }%
  }%
  \par\vspace{0.47em}%
}

\newcommand{\answerbar}[1]{%
  \vspace{0.55em}%
  \noindent\colorbox{sectiongray}{%
    \parbox{\dimexpr\textwidth-2\fboxsep\relax}{%
      \strut\textbf{#1}\strut
    }%
  }%
  \par\vspace{0.50em}%
}

\newlist{parts}{enumerate}{1}
\setlist[parts]{
  label=(\alph*),
  leftmargin=2.1em,
  itemsep=0.45em,
  topsep=0.25em,
  parsep=0pt
}

\begin{document}

% ================================================================
% PAGE 1
% ================================================================

\begin{center}
    {\Large\bfseries Math 2650 -- Quiz 04 Concept Practice}\\[0.28em]
    \small Choose the section matching the skill you want to practice.
    Answers are on the back.
\end{center}

\vspace{0.07em}

% ================================================================
% A
% ================================================================

\sectionbar{A. Position, Velocity, and the Osculating Plane}

A particle has acceleration
\[
\mathbf a(t)
=
\langle 2\cos t,\;6t+1,\;-4\rangle.
\]
At time $t=0$, its position and velocity are
\[
\mathbf r(0)=\langle0,1,-1\rangle,
\qquad
\mathbf v(0)=\langle1,2,1\rangle.
\]

\begin{parts}

\item Find the position function $\mathbf r(t)$ and determine the
particle's position at $t=1$.

\item Find an equation of the osculating plane to the particle's path
at $t=0$.

\end{parts}

\vspace{0.12em}

% ================================================================
% B
% ================================================================

\sectionbar{B. Motion with Constant Acceleration}

An object moves with constant acceleration
\[
\mathbf a(t)=\langle0,g,0\rangle.
\]
At $t=0$,
\[
\mathbf r(0)=\langle2,50,-3\rangle,
\qquad
\mathbf v(0)=\langle6,0,8\rangle.
\]
At some later time, the object is at
\[
\langle14,18,13\rangle.
\]

\begin{parts}

\item Find the later time and determine the value of $g$.

\item Find the velocity of the object at that time.

\end{parts}

\vspace{0.12em}

% ================================================================
% C
% ================================================================

\sectionbar{C. Arc Length and Distance Along a Curve}

Consider the curve
\[
\mathbf r(t)
=
\left\langle
2(\cos t+t\sin t),\;
2(\sin t-t\cos t),\;
t^2
\right\rangle,
\qquad t\geq0.
\]

\begin{parts}

\item Find the distance traveled along the curve from $t=1$ to $t=3$.

\item Starting at $\mathbf r(0)$ and moving in the direction of
increasing $t$, find the point $P$ that lies $9\sqrt2$ units along
the curve from $\mathbf r(0)$.

\end{parts}

\vspace{0.12em}

% ================================================================
% D
% ================================================================

\sectionbar{D. Tangential and Normal Acceleration and Curvature}

A particle moves along the curve
\[
\mathbf r(t)
=
\left\langle
t+\frac{t^3}{3},\;
t-\frac{t^3}{3},\;
t^2
\right\rangle,
\qquad t\geq0.
\]

\begin{parts}

\item Find the tangential component of acceleration $a_T(t)$.

\item Find the normal component of acceleration $a_N(t)$.

\item Find the curvature $\kappa(t)$.

\end{parts}

% ================================================================
% PAGE 2: ANSWER KEY
% ================================================================

\newpage

\begin{center}
    {\Large\bfseries Answer Key}\\[0.28em]
    \small Final answers only
\end{center}

\vspace{0.30em}

% ================================================================
% A ANSWERS
% ================================================================

\answerbar{A. Position, Velocity, and the Osculating Plane}

\begin{parts}

\item
\[
\boxed{
\mathbf r(t)
=
\left\langle
-2\cos t+t+2,\;
t^3+\frac12t^2+2t+1,\;
-2t^2+t-1
\right\rangle
}
\]
and
\[
\boxed{
\mathbf r(1)
=
\left\langle
3-2\cos1,\;
\frac92,\;
-2
\right\rangle
}.
\]

\item
\[
\boxed{-3x+2y-z=3}
\]

\end{parts}

\vspace{0.25em}

% ================================================================
% B ANSWERS
% ================================================================

\answerbar{B. Motion with Constant Acceleration}

\begin{parts}

\item
\[
\boxed{t=2},
\qquad
\boxed{g=-16}.
\]

\item
\[
\boxed{\mathbf v(2)=\langle6,-32,8\rangle}
\]

\end{parts}

\vspace{0.25em}

% ================================================================
% C ANSWERS
% ================================================================

\answerbar{C. Arc Length and Distance Along a Curve}

\begin{parts}

\item
\[
\boxed{8\sqrt2}
\]

\item
\[
\boxed{
P=
\left(
2(\cos3+3\sin3),\;
2(\sin3-3\cos3),\;
9
\right)
}
\]

\end{parts}

\vspace{0.25em}

% ================================================================
% D ANSWERS
% ================================================================

\answerbar{D. Tangential and Normal Acceleration and Curvature}

\begin{parts}

\item
\[
\boxed{a_T(t)=2\sqrt2\,t}
\]

\item
\[
\boxed{a_N(t)=2}
\]

\item
\[
\boxed{\kappa(t)=\frac{1}{(1+t^2)^2}}
\]

\end{parts}

\end{document}